\documentclass[10pt]{article}
\usepackage[T1]{fontenc}
\usepackage[utf8]{inputenc}
\usepackage[french]{babel}
\usepackage{amsmath,amssymb,mathrsfs,mathpazo}
\usepackage{geometry}
\geometry{a4paper,margin=2cm}
\usepackage{graphicx,array,tabularx,tabularray,colortbl,multirow,multicol,booktabs,enumitem,pifont,ulem,fancybox,diagbox,cancel}
\usepackage{pstricks,pstricks-add,pst-plot,pst-tree,pst-node,pst-text,pst-func,pst-math,pst-bspline,pst-3dplot,pst-eucl,pst-solides3d}
\let\pscalculate\undefined
\usepackage{pst-bezier}
\usepackage{tikz,pgfplots,tkz-tab}
\usetikzlibrary{arrows,arrows.meta,patterns,calc,positioning,trees,decorations.markings,shapes.geometric,shapes.misc,shapes.arrows,intersections,decorations.pathmorphing,decorations.pathreplacing,backgrounds,shadows,babel}
\pgfplotsset{compat=1.18}
\usepgfplotslibrary{fillbetween}
\usepackage[np]{numprint}
\usepackage{listings,makecell,pas-tableur,twoopt}
\let\barre\undefined
\usepackage[pstricks]{tablvar}
\definecolor{darkgreen}{RGB}{0,100,0}
\providecommand{\red}{\color{red}}
\providecommand{\blue}{\color{blue}}
\providecommand{\green}{\color{green}}
\providecommand{\black}{\color{black}}
\providecommand{\white}{\color{white}}
\providecommand{\gray}{\color{gray}}
\providecommand{\cyan}{\color{cyan}}
\providecommand{\magenta}{\color{magenta}}
\providecommand{\yellow}{\color{yellow}}
\providecommand{\violet}{\color{violet}}
\providecommand{\orange}{\color{orange}}
\providecommand{\euro}{\texteuro}
\providecommand{\EUR}{\texteuro}
\providecommand{\ieme}{\textsuperscript{e}}
\providecommand{\degres}{\ensuremath{^\circ}}
\providecommand{\degree}{\ensuremath{^\circ}}
\providecommand{\vv}[1]{\overrightarrow{#1}}
\providecommand{\Coord}[1]{#1}
\providecommand{\decosix}{}
\providecommand{\starredbullet}{\ensuremath{\bullet}}
\pagestyle{empty}
\setlength{\parindent}{0pt}
\usepackage{adjustbox}\DeclareUnicodeCharacter{2236}{:}\DeclareUnicodeCharacter{2190}{\ensuremath{\leftarrow}}
\providecommand{\R}{\mathbb{R}}
\providecommand{\N}{\mathbb{N}}
\providecommand{\Z}{\mathbb{Z}}
\providecommand{\Q}{\mathbb{Q}}
\providecommand{\D}{\mathbb{D}}
\providecommand{\C}{\mathbb{C}}
\providecommand{\e}{{\mathrm{e}}}
\providecommand{\vect}[1]{\overrightarrow{#1}}
\providecommand{\vectt}[1]{\overrightarrow{\mathrm{#1}}}
\providecommand{\vvt}[1]{\overrightarrow{\mathrm{#1}}}
\providecommand{\np}[2][]{{#2~#1}}
\providecommand{\vv}[1]{{\overrightarrow{#1}}}
\renewcommand{\d}{{\mathrm{d}}}
\providecommand{\degres}{{^\circ}}
\providecommand{\degre}{{^\circ}}
\providecommand{\blue}{\color{blue}}
\providecommand{\red}{\color{red}}
\providecommand{\enskip}{\quad}
\providecommand{\llbracket}{\lbrack\!\lbrack}
\providecommand{\rrbracket}{\rbrack\!\rbrack}
\providecommand{\euro}{{\text{€}}}
\providecommand{\no}{{\text{n°}}}
\providecommand{\interval}[2]{{[#1;#2]}}
\providecommand{\intervFO}[2]{{[#1;#2[}}
\providecommand{\dotfill}{\ldots}

\providecommand{\ds}{\displaystyle}
\providecommand{\textding}[1]{#1}
\providecommand{\Oij}{\(\mathrm{O};\vec{\imath},\vec{\jmath}\)}
\providecommand{\Oijk}{\(\mathrm{O};\vec{\imath},\vec{\jmath},\vec{k}\)}
\providecommand{\Ouv}{\(\mathrm{O};\vec{u},\vec{v}\)}
\providecommand{\intervFF}[2]{\ensuremath{[#1;#2]}}
\renewcommand{\euro}{\ifmmode\text{\texteuro}\else\texteuro\fi}
\DeclareUnicodeCharacter{00B0}{\ifmmode^{\circ}\else\textdegree\fi}
\DeclareUnicodeCharacter{20AC}{\ifmmode\text{\texteuro}\else\texteuro\fi}
% Énoncés : Éducation nationale. Transcription et corrigés : APMEP, Denis Vergès et contributeurs. Découpage et classement : CoopMaths, Sébastien Lozano, Éric Elter et contributeurs. Adaptation de la présentation : Maths Explorer.
% CC BY-SA : https://coopmaths.fr/www/about/
% https://coopmaths.fr/alea/static/bac/2022/tex/bac_2022_05_sujet1_asie_4.tex
\begin{document}
\noindent\textit{APMEP / CoopMaths. CC BY-SA. Présentation adaptée par Maths Explorer.}\par\medskip
\medskip

\emph{Principaux domaines abordés}: Étude de fonctions. Fonction logarithme.

\bigskip

Soit $f$ une fonction définie et dérivable sur $\R$. On considère les points A(1~;~3) et B(3~;~5).

On donne ci-dessous $\mathcal{C}_f$ la courbe représentative de $f$ dans un repère orthogonal du plan, ainsi que la tangente (AB) à la courbe $\mathcal{C}_f$ au point A.

\medskip

\begin{center}
\psset{unit=0.8cm,arrowsize=2pt 3}
\begin{pspicture*}(-7,-2.25)(8,6.25)
\psgrid[gridlabels=0pt,subgriddiv=1,gridwidth=0.15pt]
\psaxes[linewidth=1.25pt,labelFontSize=\scriptstyle]{->}(0,0)(-7,-2.25)(8,6.25)
\psplot[plotpoints=2000,linewidth=1.25pt,linecolor=red]{-7}{8}{x dup mul 1 add ln 3 add 2 ln sub}
\psplot[linestyle=dashed]{-7}{8}{x 2 add}
\uput[dr](1,3){A}\uput[ul](3,5){B}
\psdots[dotstyle=+,dotangle=45,dotscale=1.5](1,3)(3,5)
\uput[d](-5.5,5.75){\red $\mathcal{C}_f$}
\end{pspicture*}
\end{center}

\emph{Les trois parties de l'exercice peuvent être traitées de manière indépendante.}

\bigskip

\textbf{Partie A}

\medskip

\begin{enumerate}
\item Déterminer graphiquement les valeurs de $f(1)$ et $f'(1)$.
\item La fonction $f$ est définie par l'expression $f(x) = \ln \left(ax^2 + 1\right) + b$, où $a$ et $b$ sont des nombres réels positifs.
	\begin{enumerate}
		\item Déterminer l'expression de $f'(x)$.
		\item Déterminer les valeurs de $a$ et $b$ à l'aide des résultats précédents.

	\end{enumerate}
\end{enumerate} 

\bigskip

\textbf{Partie B}

\medskip

On admet que la fonction $f$ est définie sur $\R$ par 

\[f(x) = \ln \left(x^2 + 1\right) + 3 - \ln (2).\]

\medskip

\begin{enumerate}
\item Montrer que $f$ est une fonction paire.
\item Déterminer les limites de $f$ en $+\infty$ et en $-\infty$.
\item Déterminer l'expression de $f'(x)$.

Étudier le sens de variation de la fonction $f$ sur $\R$.

Dresser le tableau des variations de $f$ en y faisant figurer la valeur exacte du minimum ainsi que les limites de $f$ en $-\infty$ et $+\infty$.
\item À l'aide du tableau des variations de $f$, donner les valeurs du réel $k$ pour lesquelles l'équation $f(x) = k$ admet deux solutions.
\item Résoudre l'équation $f(x) = 3 + \ln 2$.
\end{enumerate}

\bigskip

\textbf{Partie C}

\medskip

On rappelle que la fonction $f$ est définie sur $R$ par $f(x) = \ln \left(x^2 + 1\right) + 3 - \ln (2)$.

\medskip

\begin{enumerate}
\item Conjecturer, par lecture graphique, les abscisses des éventuels points d'inflexion
de la courbe $\mathcal{C}_f$.
\item Montrer que, pour tout nombre réel $x$, on a : $f''(x) = \dfrac{2\left(1 - x^2\right)}{\left(x^2 + 1\right)^2}$.
\item En déduire le plus grand intervalle sur lequel la fonction $f$ est convexe.
\end{enumerate}

\end{document}
